%source mathadata
On reprend les points de référence : $A(1\,;\,2)$ pour le chiffre 1, $B(5\,;\,2)$ pour le chiffre 3 et $C(3\,;\,6)$ pour le chiffre 8.

		\begin{multicols}{2}

			\begin{MADenumerate}

				\item Déterminer l'équation de la droite qui sépare les zones des chiffres 1 et 3.

				\begin{exemple}[Aide]
					Soit $M(x\,;\,y)$ un point sur la droite qui sépare les zones de $A(1\,;\,2)$ et $B(5\,;\,2)$.
					On a donc :
					\begingroup\small
					\begin{align*}
						MA                            & \; \ldots\; MB                        \\
						MA^2                          & \; \ldots\; MB^2                      \\
						(x-\ldots)^2+(y-\ldots)^2     & \; \ldots\; (x-\ldots)^2+(y-\ldots)^2 \\
						\text{Après simplification :} &                                       \\
						y                             & = \ldots x+\ldots
					\end{align*}
					\endgroup
				\end{exemple}

				\item Déterminer de même les équations des droites séparant les zones des chiffres 1 et 8, puis des chiffres 3 et 8.

				\item Pour chaque droite, chercher deux points avant de tracer. Tracer les trois droites dans le repère ci-contre et mettre en évidence les trois zones d'estimation.

				\item Que remarquez-vous pour les trois droites obtenues ?

				\item On donne le point $P(3\,;\,3{,}5)$. Vérifier que ce point vérifie les trois équations trouvées. Pouviez-vous prévoir ce résultat ?

				\item (Bonus) Que représentent ces trois droites pour le triangle $ABC$ ? Comment appelle-t-on le point $P$ ?

				\item (Bonus) Vérifier que $P$ est à égale distance de $A$, $B$ et $C$. Que cela signifie-t-il géométriquement ?

			\end{MADenumerate}

			\begin{figure}[H]
				\centering
				\begin{tikzpicture}[line cap=round,line join=round,>=triangle 45,x=1cm,y=1cm,scale=0.7]
					\begin{axis}[
							x=1cm,y=1cm,
							axis lines=middle,
							ymajorgrids=true,
							xmajorgrids=true,
							xmin=-1,
							xmax=8,
							ymin=-1,
							ymax=8,
							xtick={0,1,...,7},
							ytick={0,1,...,7},]
					\end{axis}
				\end{tikzpicture}
			\end{figure}

		\end{multicols}
